% Generated draft: requires mathematical and source review.
\documentclass[11pt]{article}
\usepackage[margin=25mm]{geometry}
\usepackage{fontspec}
\setmainfont{Latin Modern Roman}
\usepackage{amsmath,amssymb,mathtools,unicode-math}
\setmathfont{Latin Modern Math}
\usepackage{enumitem,needspace,xurl,hyperref}
\setlength{\parindent}{0pt}
\setlength{\parskip}{5pt}
\setlength{\emergencystretch}{3em}
\newcommand{\N}{\mathbb{N}}
\newcommand{\Z}{\mathbb{Z}}
\newcommand{\Q}{\mathbb{Q}}
\newcommand{\R}{\mathbb{R}}
\newcommand{\C}{\mathbb{C}}
\newcommand{\F}{\mathbb{F}}
\newcommand{\A}{\mathbb{A}}
\newcommand{\PP}{\mathbb{P}}
\newcommand{\E}{\mathbb{E}}
\newcommand{\eps}{\varepsilon}
\newcommand{\dd}{\,\mathrm{d}}
\begin{document}
\section*{UnsolvedMath ID: 30003809}
\textbf{Source problem code:} catalogue entry
\subsection{Definitions and mathematical statement}
The supplied problem concerns the operators
\[
A_t=\lambda_t^{*}+\lambda_t+\rho_t^{*}+\rho_t,
\]
which the source describes as self-adjoint bimonotone limit operators. Here \(^{*}\) denotes the Hilbert-space adjoint. The record does not define \(\lambda_t\) or \(\rho_t\), their underlying Hilbert space, their domains, the range of \(t\), or the probability state used to take their distributions. These data must be recovered from the source before the target measures are uniquely specified.

To make the meaning of distribution precise conditionally, suppose that a self-adjoint realization of \(A_t\) and a probability state \(\varphi_t\) on its bounded Borel functional calculus have been specified, with \(\varphi_t\) countably additive on spectral projections. Writing \(E_{A_t}\) for the spectral measure of that realization, its distribution is the Borel probability measure \(\mu_t\) on \(\R\) defined by
\[
\mu_t(B)=\varphi_t\bigl(E_{A_t}(B)\bigr)
\qquad\text{for every Borel set }B\subseteq\R.
\]
Equivalently,
\[
\int_{\R} f(x)\,\dd\mu_t(x)=\varphi_t\bigl(f(A_t)\bigr)
\]
for every bounded Borel function \(f:\R\to\C\). This describes the general meaning of distribution; it does not supply the missing bimonotone construction or select a state.

\textbf{Requested problem.} Give a concrete description of the measures \(\mu_t\) associated with the source's self-adjoint bimonotone limit operators \(A_t\), for every parameter \(t\) allowed by that construction. The supplied record does not specify what counts as a concrete description, and does not establish that these measures have densities. The parameter range, operator construction, state, and any domain or closure conventions remain unresolved qualifications of this statement.
\subsection{Short English statement}
Describe explicitly the real probability laws of the self-adjoint bimonotone limit operators \(\lambda_t^{*}+\lambda_t+\rho_t^{*}+\rho_t\). The supplied record leaves the operators, parameter range, and probability state undefined, so it does not yet determine a unique family of laws. Its open-status label is supplied curation, not an independently verified conclusion.
\subsection{Sources}
\begin{itemize}
\item Local UnsolvedMath registry, permanent ID \(30003809\), title \emph{Distributions of Self-Adjoint Bimonotone Limit Operators}, source problem code \texttt{catalogue entry}, in the Oberwolfach Reports Open Problems collection. The record gives proposed year 2018 and names no individual proposer.
\item Supplied source citation: \emph{Interactions between Operator Space Theory and Quantum Probability with Applications to Quantum Information} (2019), Oberwolfach Reports, pp.\ 1293--1379, DOI \texttt{10.4171/owr/2018/22}. \url{https://doi.org/10.4171/owr/2018/22}. The background describes these reports as documenting mathematical workshops and their open-problem discussions.
\item Supplied literature source: \emph{Bi monotone Brownian motion}. \url{https://arxiv.org/abs/1708.03510}.
\item Supplied literature source: \emph{Published bi monotone Brownian motion}. \url{https://doi.org/10.1142/9789811275999_0005}.
\item Literature assessment dated exactly \texttt{2026-08-21}: the supplied curation labels the problem open and states that Gerhold's later publication develops bi monotone Brownian motion and its laws but does not give the requested concrete description of the corresponding probability measures. This is supplied dated literature triage derived from the OWR curation, not an independently verified assessment. Its time-sensitive status and cited claims require independent verification.
\end{itemize}
\end{document}
